\begin{proposition}[The root base and its explicit cohomology]\hypertarget{tag-01SM}{}\label{inv:D1.1}
\StatementRecord{inv:D1.1}{proved}{01SM}{the definitions and constructions stated in the proof}

Let $\rho:\mathcal B\to B$ be the analytic root curve with orders $3$ and
$4$ at $p_1,p_2$. Write $R_1,R_2$ for its tautological root line bundles.
Its Picard group is infinite cyclic, with generator
$R=R_1\otimes R_2^{-1}$, and
\[
 R_1=R^4,\quad R_2=R^3,\quad \rho^*\OO_B(1)=R^{12},
 \quad K_{\mathcal B}=R^{-7}.
\]
We abbreviate $R^n$ to $\OO_{\mathcal B}(n)$. Choose sections
$x\in H^0(R^3)$ and $y\in H^0(R^4)$ vanishing at $p_2$ and $p_1$,
respectively. For every integer $n$, put
\[
 r_1(n)\equiv n\pmod3,\quad 0\le r_1(n)<3,
 \qquad r_2(n)\equiv3n\pmod4,\quad0\le r_2(n)<4,
\]
\[
 d(n)=\frac{n-4r_1(n)-3r_2(n)}{12}.
\]
Then
\begin{equation}
 \hypertarget{tag-01P3}{}\label{fam:eq:ab-root-degree}
\rho_*\OO_{\mathcal B}(n)=\OO_B(d(n)),\qquad
 H^i(\mathcal B,\OO(n))=H^i(B,\OO(d(n))).                 
\end{equation}
Equivalently, $h^0(\OO(n))$ is the number of pairs
$(a,b)\in\Z_{\ge0}^2$ with $3a+4b=n$, and
$h^1(\OO(n))=h^0(\OO(-n-7))$.
\end{proposition}



\begin{proof}
The Picard group and graded cohomology are those of
\cite[Lemma~6.1, p.~74]{DuEtAl2026v1}; the following chart maps specify
the coarse invariant direct image used here.
\begin{enumerate}[label=\arabic*.,leftmargin=*]
\item On the root charts the group acts on a disc coordinate by a primitive
$m_j$th root of unity. The stabilizer character of any line bundle is
removed uniquely by tensoring with $R_1^{-r_1}R_2^{-r_2}$, with the
exponents in the displayed ranges. A line bundle with trivial stabilizer
characters descends: an equivariant nonzero local frame can be averaged,
and its invariant frame gives the descent transition function. It is
therefore the pullback of a unique line bundle on $B$. Consequently
\[
 \operatorname{Pic}(\mathcal B)=
 \langle H,R_1,R_2\mid 3R_1=H,\;4R_2=H\rangle=\Z R.
\]
The weighted quotient presentation is
\[
 \mathcal B\simeq[(\mathbb C^2\setminus\{0\})/\mathbb C^*],
 \quad\lambda(x,y)=(\lambda^3x,\lambda^4y),\qquad
 t=\frac{y^3}{y^3-x^4}.
\]
The root sections are $(\mathcal O(4),y)$ and $(\mathcal O(3),x)$.
They give isomorphisms on each root chart, and the same coarse
coordinate on their ordinary overlap makes those isomorphisms glue.
Additive notation in the Picard presentation means tensor product. The determinant of the coordinate change $t=s^m$ gives
$K_{\mathcal B}=\rho^*K_B\otimes R_1^2R_2^3=R^{-24+8+9}$.

\item An invariant section of $R_j^{r_j}$ on its chart is $s^{r_j}$ times a
holomorphic function of $t=s^{m_j}$. It follows, by the local invariant
comparison on overlaps, that
$\rho_*(\rho^*\OO(k)R_1^{r_1}R_2^{r_2})=\OO(k)$.
In characteristic zero taking finite-group invariants is exact: the
averaging idempotent splits the invariant submodule. Thus there are no
higher direct images for $\rho$, and \eqref{fam:eq:ab-root-degree} follows. This proves the
cohomology comparison with its canonical invariant map. Finally every monomial of degree $n$ is
$x^{r_2}y^{r_1}$ times a homogeneous polynomial of degree $d(n)$ in
$x^4,y^3$. This gives the stated count of sections. The negative-degree
\v Cech complex on $x\ne0$, $y\ne0$ has cokernel represented by
$x^{-a}y^{-b}$ with $a,b\ge1$ and $3a+4b=-n$; replacing $(a,b)$ by
$(a-1,b-1)$ proves the $h^1$ formula. \end{enumerate}
\end{proof}